% =========================================================
% Modulus: Algebra of Absolute Value
% =========================================================

\section{Modulus}
\label{sec:modulus}
% ---------------------------------------------------------
% def:absolute-value
% ---------------------------------------------------------

\begin{definitionbox}{Definition (Absolute Value)}
\begin{definition}[Absolute Value]
\label{def:absolute-value}
Let $a \in \mathbb{R}$.  The \emph{absolute value} (or \emph{modulus}) of
$a$ is
\[
  \lvert a \rvert
  \;:=\;
  \begin{cases}
    a  & \text{if } a \ge 0, \\
    -a & \text{if } a < 0.
  \end{cases}
\]
Equivalently, $\lvert a \rvert := \max(a, -a)$.
\end{definition}
\end{definitionbox}

\begin{remark*}[Standard quantified statement]
\[
  \lvert a \rvert
  \;=\;
  \max(a,\,-a)
  \;\coloneqq\;
  \begin{cases}
    a  & \text{if } a \ge 0, \\
    -a & \text{if } a < 0.
  \end{cases}
\]
\end{remark*}

\begin{remark*}[Interpretation]
The absolute value measures magnitude independent of sign.  On the real
line it equals the distance from $a$ to $0$: $\lvert a \rvert = d(a, 0)$
under the standard metric.  The two-branch definition partitions
$\mathbb{R}$ at $0$ and reflects negative values to positive ones.
\end{remark*}

\begin{remark*}[Dependencies]
\hyperref[def:real-ordered-field]{Ordered field},
\hyperref[thm:q-is-ordered-field]{$\mathbb{R}$ is an ordered field}.
\end{remark*}

% ---------------------------------------------------------
% thm:absolute-value-nonneg
% ---------------------------------------------------------

\begin{theorem}[Absolute Value Is Nonnegative]
\label{thm:absolute-value-nonneg}
For all $a \in \mathbb{R}$,
\[
  \lvert a \rvert \ge 0.
\]
\end{theorem}



\begin{remark*}[Proof]
\hyperref[prf:absolute-value-nonneg]{Go to proof.}
\end{remark*}

\begin{remark*}[Standard quantified statement]
\[
  \forall a \in \mathbb{R},\quad \lvert a \rvert \ge 0.
\]
\end{remark*}

\begin{remark*}[Dependencies]
\hyperref[def:absolute-value]{Absolute value},
\hyperref[thm:squares-nonnegative]{squares are nonnegative},
\hyperref[def:real-ordered-field]{ordered field}.
\end{remark*}

% ---------------------------------------------------------
% thm:absolute-value-zero-iff-zero
% ---------------------------------------------------------

\begin{theorem}[Absolute Value Zero Iff Zero]
\label{thm:absolute-value-zero-iff-zero}
For all $a \in \mathbb{R}$,
\[
  \lvert a \rvert = 0 \iff a = 0.
\]
\end{theorem}



\begin{remark*}[Proof]
\hyperref[prf:absolute-value-zero-iff-zero]{Go to proof.}
\end{remark*}

\begin{remark*}[Standard quantified statement]
\[
  \forall a \in \mathbb{R},\quad
  \lvert a \rvert = 0 \iff a = 0.
\]
\end{remark*}

\begin{remark*}[Dependencies]
\hyperref[def:absolute-value]{Absolute value},
\hyperref[thm:absolute-value-nonneg]{absolute value is nonnegative},
\hyperref[def:real-ordered-field]{ordered field}.
\end{remark*}

% ---------------------------------------------------------
% thm:absolute-value-self-or-neg
% ---------------------------------------------------------

\begin{theorem}[Absolute Value Equals Self or Negation]
\label{thm:absolute-value-self-or-neg}
For all $a \in \mathbb{R}$, either $\lvert a \rvert = a$ or
$\lvert a \rvert = -a$.
\end{theorem}



\begin{remark*}[Proof]
\hyperref[prf:absolute-value-self-or-neg]{Go to proof.}
\end{remark*}

\begin{remark*}[Standard quantified statement]
\[
  \forall a \in \mathbb{R},\quad
  \lvert a \rvert = a \;\lor\; \lvert a \rvert = -a.
\]
\end{remark*}

\begin{remark*}[Dependencies]
\hyperref[def:absolute-value]{Absolute value}.
\end{remark*}

% ---------------------------------------------------------
% thm:absolute-value-symmetric
% ---------------------------------------------------------

\begin{theorem}[Absolute Value Is Symmetric]
\label{thm:absolute-value-symmetric}
For all $a \in \mathbb{R}$,
\[
  \lvert -a \rvert = \lvert a \rvert.
\]
\end{theorem}



\begin{remark*}[Proof]
\hyperref[prf:absolute-value-symmetric]{Go to proof.}
\end{remark*}

\begin{remark*}[Standard quantified statement]
\[
  \forall a \in \mathbb{R},\quad \lvert -a \rvert = \lvert a \rvert.
\]
\end{remark*}

\begin{remark*}[Dependencies]
\hyperref[def:absolute-value]{Absolute value},
\hyperref[thm:double-negation]{double negation}.
\end{remark*}

% ---------------------------------------------------------
% thm:absolute-value-product
% ---------------------------------------------------------

\begin{theorem}[Absolute Value of a Product]
\label{thm:absolute-value-product}
For all $a, b \in \mathbb{R}$,
\[
  \lvert ab \rvert = \lvert a \rvert\,\lvert b \rvert.
\]
\end{theorem}



\begin{remark*}[Proof]
\hyperref[prf:absolute-value-product]{Go to proof.}
\end{remark*}

\begin{remark*}[Standard quantified statement]
\[
  \forall a, b \in \mathbb{R},\quad
  \lvert ab \rvert = \lvert a \rvert\,\lvert b \rvert.
\]
\end{remark*}

\begin{remark*}[Dependencies]
\hyperref[def:absolute-value]{Absolute value},
\hyperref[thm:negation-of-product]{negation of product},
\hyperref[cor:product-of-negatives]{product of negatives},
\hyperref[thm:squares-nonnegative]{squares are nonnegative}.
\end{remark*}

% ---------------------------------------------------------
% thm:absolute-value-quotient
% ---------------------------------------------------------

\begin{theorem}[Absolute Value of a Quotient]
\label{thm:absolute-value-quotient}
For all $a \in \mathbb{R}$ and $b \in \mathbb{R}$ with $b \ne 0$,
\[
  \left\lvert \frac{a}{b} \right\rvert
  = \frac{\lvert a \rvert}{\lvert b \rvert}.
\]
\end{theorem}



\begin{remark*}[Proof]
\hyperref[prf:absolute-value-quotient]{Go to proof.}
\end{remark*}

\begin{remark*}[Standard quantified statement]
\[
  \forall a \in \mathbb{R},\;
  \forall b \in \mathbb{R} \setminus \{0\},\quad
  \left\lvert \frac{a}{b} \right\rvert
  = \frac{\lvert a \rvert}{\lvert b \rvert}.
\]
\end{remark*}

\begin{remark*}[Dependencies]
\hyperref[def:absolute-value]{Absolute value},
\hyperref[thm:absolute-value-product]{absolute value of a product},
\hyperref[thm:absolute-value-zero-iff-zero]{absolute value zero iff zero},
\hyperref[thm:uniqueness-of-multiplicative-inverse]{uniqueness of
multiplicative inverse}.
\end{remark*}

% ---------------------------------------------------------
% thm:absolute-value-bounds
% ---------------------------------------------------------

\begin{theorem}[Bound by Modulus]
\label{thm:absolute-value-bounds}
For all $a \in \mathbb{R}$,
\[
  -\lvert a \rvert \le a \le \lvert a \rvert.
\]
\end{theorem}



\begin{remark*}[Proof]
\hyperref[prf:absolute-value-bounds]{Go to proof.}
\end{remark*}

\begin{remark*}[Standard quantified statement]
\[
  \forall a \in \mathbb{R},\quad
  -\lvert a \rvert \le a \le \lvert a \rvert.
\]
\end{remark*}

\begin{remark*}[Dependencies]
\hyperref[def:absolute-value]{Absolute value},
\hyperref[thm:absolute-value-nonneg]{absolute value is nonnegative},
\hyperref[def:real-ordered-field]{ordered field}.
\end{remark*}

% ---------------------------------------------------------
% thm:absolute-value-le-iff
% ---------------------------------------------------------

\begin{theorembox}{Theorem (Nonstrict Modulus Inequality)}
\begin{theorem}[Nonstrict Modulus Inequality]
\label{thm:absolute-value-le-iff}
For all $a \in \mathbb{R}$ and $r \ge 0$,
\[
  \lvert a \rvert \le r \iff -r \le a \le r.
\]
\end{theorem}
\end{theorembox}



\begin{remark*}[Proof]
\hyperref[prf:absolute-value-le-iff]{Go to proof.}
\end{remark*}

\begin{remark*}[Standard quantified statement]
\[
  \forall a \in \mathbb{R},\;
  \forall r \in \mathbb{R},\;
  r \ge 0 \;\Rightarrow\;
  \bigl(\lvert a \rvert \le r \iff -r \le a \le r\bigr).
\]
\end{remark*}

\begin{remark*}[Interpretation]
This biconditional is the primary tool for replacing a modulus inequality
with a pair of simultaneous linear inequalities.  It underlies the
$\varepsilon$-$\delta$ manipulations throughout analysis: the condition
$\lvert x - a \rvert \le \varepsilon$ is equivalent to
$a - \varepsilon \le x \le a + \varepsilon$.
\end{remark*}

\begin{remark*}[Dependencies]
\hyperref[def:absolute-value]{Absolute value},
\hyperref[thm:absolute-value-bounds]{bound by modulus},
\hyperref[def:real-ordered-field]{ordered field}.
\end{remark*}

% ---------------------------------------------------------
% thm:absolute-value-lt-iff
% ---------------------------------------------------------

\begin{theorem}[Strict Modulus Inequality]
\label{thm:absolute-value-lt-iff}
For all $a \in \mathbb{R}$ and $r > 0$,
\[
  \lvert a \rvert < r \iff -r < a < r.
\]
\end{theorem}



\begin{remark*}[Proof]
\hyperref[prf:absolute-value-lt-iff]{Go to proof.}
\end{remark*}

\begin{remark*}[Standard quantified statement]
\[
  \forall a \in \mathbb{R},\;
  \forall r \in \mathbb{R},\;
  r > 0 \;\Rightarrow\;
  \bigl(\lvert a \rvert < r \iff -r < a < r\bigr).
\]
\end{remark*}

\begin{remark*}[Dependencies]
\hyperref[def:absolute-value]{Absolute value},
\hyperref[thm:absolute-value-le-iff]{nonstrict modulus inequality},
\hyperref[def:real-ordered-field]{ordered field}.
\end{remark*}

% ---------------------------------------------------------
% thm:triangle-inequality
% ---------------------------------------------------------

\begin{theorembox}{Theorem (Triangle Inequality)}
\begin{theorem}[Triangle Inequality]
\label{thm:triangle-inequality}
For all $a, b \in \mathbb{R}$,
\[
  \lvert a + b \rvert \le \lvert a \rvert + \lvert b \rvert.
\]
\end{theorem}
\end{theorembox}



\begin{remark*}[Proof]
\hyperref[prf:triangle-inequality]{Go to proof.}
\end{remark*}

\begin{remark*}[Standard quantified statement]
\[
  \forall a, b \in \mathbb{R},\quad
  \lvert a + b \rvert \le \lvert a \rvert + \lvert b \rvert.
\]
\end{remark*}

\begin{remark*}[Interpretation]
The triangle inequality is the foundational estimate of analysis.  Every
convergence argument, continuity proof, and metric-space result that
requires bounding a sum reduces, at some step, to this inequality.  The
name reflects its geometric meaning: the length of one side of a triangle
does not exceed the sum of the other two.
\end{remark*}

\begin{remark*}[Dependencies]
\hyperref[def:absolute-value]{Absolute value},
\hyperref[thm:absolute-value-bounds]{bound by modulus},
\hyperref[thm:absolute-value-le-iff]{nonstrict modulus inequality},
\hyperref[def:real-ordered-field]{ordered field}.
\end{remark*}

% ---------------------------------------------------------
% thm:reverse-triangle-inequality
% ---------------------------------------------------------

\begin{theorem}[Reverse Triangle Inequality]
\label{thm:reverse-triangle-inequality}
For all $a, b \in \mathbb{R}$,
\[
  \lvert\, \lvert a \rvert - \lvert b \rvert \,\rvert
  \le \lvert a - b \rvert.
\]
\end{theorem}



\begin{remark*}[Proof]
\hyperref[prf:reverse-triangle-inequality]{Go to proof.}
\end{remark*}

\begin{remark*}[Standard quantified statement]
\[
  \forall a, b \in \mathbb{R},\quad
  \lvert\, \lvert a \rvert - \lvert b \rvert \,\rvert
  \le \lvert a - b \rvert.
\]
\end{remark*}

\begin{remark*}[Interpretation]
The reverse triangle inequality provides a lower bound on $\lvert a - b
\rvert$ in terms of the difference of moduli.  It is used to bound the
modulus of a difference from below, particularly when proving that a
function is not uniformly continuous or that a sequence does not converge.
\end{remark*}

\begin{remark*}[Dependencies]
\hyperref[def:absolute-value]{Absolute value},
\hyperref[thm:triangle-inequality]{triangle inequality},
\hyperref[thm:absolute-value-symmetric]{absolute value is symmetric}.
\end{remark*}

% ---------------------------------------------------------
% thm:absolute-value-sum-bound
% ---------------------------------------------------------

\begin{theorem}[Generalised Triangle Inequality]
\label{thm:absolute-value-sum-bound}
For all $n \in \mathbb{N}$ and $a_1, \ldots, a_n \in \mathbb{R}$,
\[
  \left\lvert \sum_{k=1}^{n} a_k \right\rvert
  \le \sum_{k=1}^{n} \lvert a_k \rvert.
\]
\end{theorem}



\begin{remark*}[Proof]
\hyperref[prf:absolute-value-sum-bound]{Go to proof.}
\end{remark*}

\begin{remark*}[Standard quantified statement]
\[
  \forall n \in \mathbb{N},\;
  \forall a_1, \ldots, a_n \in \mathbb{R},\quad
  \left\lvert \sum_{k=1}^{n} a_k \right\rvert
  \le \sum_{k=1}^{n} \lvert a_k \rvert.
\]
\end{remark*}

\begin{remark*}[Dependencies]
\hyperref[def:absolute-value]{Absolute value},
\hyperref[thm:triangle-inequality]{triangle inequality}.
\end{remark*}
